\documentclass[11pt,a4paper]{article}

% --- Language and Font Setup ---
\usepackage[english,greek]{babel}
\usepackage{fontspec}

% Overleaf usually has 'Times New Roman' or 'FreeSerif'
\setmainfont{FreeSerif} 

% --- Math and Logic ---
\usepackage{amsmath, amssymb, amsthm}
\usepackage{listings}
\usepackage{xcolor}

% --- LOOP Language Definition ---
\lstdefinelanguage{LOOP}{
    keywords={LOOP, DO, END, IF, THEN, ELSE, :=},
    sensitive=true,
    morecomment=[l]{//},
    morestring=[b]"
}

\lstset{
    language=LOOP,
    basicstyle=\ttfamily\small,
    keywordstyle=\color{blue}\bfseries,
    commentstyle=\color{gray}\itshape,
    frame=single,
    tabsize=2,
    breaklines=true,
    numbers=left,
    numberstyle=\tiny\color{gray}
}

% --- Document Info ---
\title{Computational Models: LOOP Excersises}
\author{Vasilis Ioannou}
\date{\today}

\begin{document}

\maketitle

\section*{LOOP Models}

LOOP models represent a formal programming language designed to characterize the class of \textbf{primitive recursive functions}. Unlike General Recursive functions that may run indefinitely, LOOP models are intentionally restricted to ensure that every program eventually terminates. 

The core philosophy behind the LOOP model is that the number of iterations in any loop must be determined \textit{before} the loop begins its execution, preventing the possibility of infinite recursion or ``hanging'' states.

\subsection*{Formal Definition}

\begin{mybox}{Syntax of LOOP Programs}
A LOOP program is defined inductively over a set of variables $x_1, x_2, \dots, x_n$, which hold non-negative integers $\mathbb{N}_0$. The syntax consists of:
    \item \textbf{Base Commands:}
    \begin{itemize}
        \item \textit{Zeroing:} $x_i \leftarrow 0$ (Assigns the value zero to variable $x_i$).
        \item \textit{Successor:} $x_i \leftarrow x_j + 1$ (Increments $x_j$ and stores it in $x_i$).
    \end{itemize}
    \item \textbf{Composition:} If $P_1$ and $P_2$ are LOOP programs, then $P_1; P_2$ is a LOOP program (sequential execution).
    \item \textbf{Iteration:} If $P$ is a LOOP program, then the construct 
    \[ \texttt{LOOP } x_i \texttt{ DO } P \texttt{ END} \]
    is a LOOP program. The command $P$ is executed exactly $n$ times, where $n$ is the value of $x_i$ at the moment the loop starts. 
\end{mybox}

\subsection*{The ``While'' Gap and the Ackermann Function}

The LOOP model is strictly less powerful than the \textbf{WHILE model} (or Turing machines). While every LOOP program is computable, not every computable function can be expressed as a LOOP program. 

The most famous counter-example is the \textbf{Ackermann function}, denoted as $A(m, n)$. It is a total computable function that is not primitive recursive because its growth rate is so extreme that it eventually exceeds the capacity of any fixed-nesting LOOP program.

\subsection*{The Ackermann Function Definition}
For $m, n \geq 0$, the function $A(m, n)$ is defined as:
\[
A(m, n) = 
\begin{cases} 
n + 1 & \text{if } m = 0 \\
A(m-1, 1) & \text{if } m > 0 \text{ and } n = 0 \\
A(m-1, A(m, n-1)) & \text{if } m > 0 \text{ and } n > 0 
\end{cases}
\]

Because the Ackermann function uses \textit{nested recursion} (where the function calls itself as an argument to itself), it requires a \texttt{WHILE} loop logic that isn't bound by a pre-determined iteration count.

\subsection*{Assign function}
Below I assigned variables in the form : x:=y for simplicity
This is done by the following function:
\begin{lstlisting}

x := 0;
LOOP i in 1 to y DO
    x := x + 1
END

\end{lstlisting}

\pagebreak

\section*{Excersise 1}

\subsection*{α) Multiplication $\text{mult}(x, y)$}
The Multiplication is implemented with two nested Loops.

\begin{lstlisting}
// input: x, y | output: z = x*y
res := 0
LOOP x in 1 to y DO
   LOOP z in 1 to x DO
      res := res + 1
   END
END
\end{lstlisting}

\subsection*{β)  Subtraction $\text{sub}(x,y)$}
Modeling the subtraction using for loops is a but trickier since we can only increment our variables. 
The main idea here is that we can essentially do the operation x := x-1 by using a variable that stores the previous value of our first variable.

\begin{lstlisting}
// input: x,y | output: target = x - y

prev:= x
res := 0;
LOOP j in 1 to y DO
    // Pref(x) returns x-1
    pred := 0
    LOOP i in 2 to prev DO
        pred := pred + 1
    END
    prev := pred
END
res := prev

\end{lstlisting}

\subsection*{β) If - zero - then $\text{ifzero}(x)$}
This is a function that returns 1 if x == 0 else returns 0.
The main idea here is that the for loop: 
    LOOP i in x to b DO
        ...
    END
Only enters in the section of the loop if  x < b.

Here was also assume x >= 0  

\begin{lstlisting}
// input: x | output: res = 1 when x == 0 else 0 

res := 0
LOOP i in x to 1 DO
    res := 1
END

\end{lstlisting}

\section*{Excersise 2}
Here we write a program in LOOP modeling for:
f(x) := if g(x) = 0 then h1(x) else h2(x)

Here g(x),h1(x),h2(x) are other LOOP programs

\subsection*{If g(x) = 0 then h1(x) else h2(x)}
\begin{lstlisting}

// input x, g(x) = x | output res

gx_res = g(x)
do_then := 1 // Here 1 means we skip it
do_else := 0

LOOP i in g(x) to 1 DO
    do_then := 0
    do_else := 1
END

LOOP i in do_then to 1 DO
    res := h1(x)
END

LOOP i in do_else to 1 DO
    res := h2(x)
END
\end{lstlisting}

\pagebreak

\section*{Excersise 3}
Here we have to implement


div(x,y) = x div y for y > 0 else 0
\subsection*{div(x,y))}
\begin{lstlisting}

quotient := 0
remainder := x
terminate := 0
LOOP i in y to 1 DO
    terminate := 1
    quotient :=0
END

LOOP k in terminate to 1 DO
    LOOP i in 1 to x DO
        check := sub(y, remainder);
        can_subtract := iszero(check);

        LOOP j in 1 to can_substract DO
            quotient := quotient + 1;
            remainder := sub(remainder, y);
        END
    END
END

\end{lstlisting}

\pagebreak

\section*{Excersise 4}
Here we have to implement pow(x,y), factorial(x,y), comb(n,a) and sqrt(x)
\subsection*{div(x,y))}
\begin{lstlisting}
// input x, y | output prod

prod:=x
LOOP i in 1 to y DO
    prod := mult(prod,x);
END

\end{lstlisting}

\subsection*{factorial(x,y))}
\begin{lstlisting}
// input x | output fact

fact:=x
temp:=sub(x,1)
LOOP i in 2 to x DO
    prod := mult(prod,temp);
    temp := sub(temp,1);
END

\end{lstlisting}

\subsection*{comb(x,y))}

comb(x,y) = x! / (y! * (x-y)!)

As we can see when we buld our basic functions LOOP modeling looks like every other programming language

\begin{lstlisting}
// input x, y | output comb

num   := fact(x)
denum := fact(y)
temp  := sub(x,y)
temp  := fact(temp)
denum := mult(denum,temp)

comb := div(num,denum);

\end{lstlisting}

\pagebreak

\subsection*{sqrt(x)}
To simulate the square root in LOOP modeling, you need to find the integer square root, 
which is the largest integer i such that i^2≤x.

To handle the "if" condition using only your basic functions, we use the property of monus (limited subtraction):

sub(1, sub(a, b)) returns 1 if a≤b
sub(1, sub(a, b)) returns 0 if a>b

\begin{lstlisting}
// input x | output res

// Input: x
// Output: result

result = 0;
current_i = 1;

LOOP i in 1 to x DO
    square = mult(current_i, current_i);
    
    // Check if square <= x
    // Logic: 1 - (square - x)
    is_le = sub(1, sub(square, x));
    
    // If it was <= x, is_le is 1, so we add it to the result
    result = add(result, is_le);
    
    // Increment our counter for the next iteration
    current_i = add(current_i, 1);
END

\end{lstlisting}
\end{document}




